\section{Introduction}

Large language model agents can repair real software bugs: they read a
repository, edit code, run tests, and recover from failures. Every one of those
attempts---successful or not---produces information about how to localize, edit,
validate, and recover on the next task. Almost all of it is thrown away the
moment the benchmark scores the trial. Recycling it into external
\emph{skills}---compact textual procedures injected into future prompts---is an
attractive way to make a frozen agent better without retraining its weights.

The attraction hides a trap. A raw trajectory is full of things that must not
leak forward: concrete file paths, the exact shape of a gold patch, transient
environment errors, and incidental identifiers of the benchmark instance. Inject
those into a future prompt and the system improves by \emph{memorization}, not
transfer---and on a held-out benchmark this is indistinguishable from progress
until someone audits the prompt. The opposite failure is just as common: a
one-shot ``summarize the trajectory into advice'' step yields instructions like
``run focused tests'' that are too generic to change behavior, or over-specific
rules that fire on the wrong task and cause regressions. We argue that most
trace-to-skill pipelines sit on one of these two failure modes and report the
resulting numbers as if they measured transfer.

The question we take seriously is therefore: \emph{can a coding-agent system
learn reusable software-engineering skills from training-time verifier feedback
without using test-time verifier feedback and without leaking task-specific
solutions?} Answering it requires more than a better summarizer. It requires (i)
a disciplined notion of what is allowed to become a transferable skill, and (ii)
a hard, auditable boundary on when the verifier may influence memory.

We propose \method{}, \textbf{hierarchical skill evolution}. The central design
commitment is restraint: \emph{not every trace artifact is a skill}. A task
attempt first becomes \emph{task evidence}---a structured record of public
signals (touched and edited paths, test commands, diagnostics, selected skills,
failure signatures, verifier outcome). Task evidence is never rendered as a
deployed skill. A Skill Writer drafts a repo-level skill only from \emph{multiple}
evidence records in the \emph{same} repository, and a failure-mode skill only
from patterns that \emph{recur across} repositories---localization drift, weak
validation, no-diff exits, timeout recovery. A Skill Evaluator then scores each
candidate for utility, support, and risk and decides whether to accept, reject,
revise, or retain it as memory-only evidence. Every upward step is gated; a
single lucky trace cannot mint a global skill.

The boundary is the second commitment, and it is hard rather than aspirational.
SWEGym verifiers are available at training time, so we use them to label
outcomes, calibrate the evaluator, and gate candidate states on a repo-isolated
validation split. On SWE-bench Verified and other downstream benchmarks, no
official verifier label, hidden test result, gold patch, or future task outcome
may influence a skill update---full stop. This is what makes a transfer claim on
those benchmarks meaningful rather than a restatement of leaked supervision.

\paragraph{Contributions.}
\begin{itemize}
    \item We recast software-agent skill evolution as a hierarchical,
    evidence-gated process that cleanly separates task evidence, repo-level
    skills, cross-repo failure-mode skills, and the writer/evaluator policies
    that produce and police them.
    \item We define a verifier-boundary protocol that licenses training-time use
    of SWEGym verifier outcomes while forbidding any verifier signal on the
    test-time update path, turning leakage from a reviewer's worry into a
    structural property of the system.
    \item We instantiate the framework end to end for the Pi coding-agent
    harness with an OpenAI-compatible model interface, Harbor/E2B execution,
    versioned skill archives, promotion logs, validation gates, and a
    deterministic evaluator shell that can be swapped for a learned policy
    without changing the loop contract.
    \item We specify an evaluation protocol over a GLM-5.2/GPT-5.5
    $\times$ Pi/Claude-style matrix that separates within-cell evolution, model
    transfer, and harness transfer across SWE-bench Verified, DeepSWE, and an
    out-of-domain Terminal-Bench 2.0/2.1 extrapolation, treating exact-task
    memory only as an explicitly labeled adaptation upper bound.
\end{itemize}
